%% Lecture 07 — Electromagnetic Lagrangian, Lorentz Force, and Gauge Invariance
\section{Lecture 07 — Electromagnetic Lagrangian, Lorentz Force, and Gauge Invariance}\label{sec:lec07}

\begin{center}
\begin{tabular}{ll}
  \textbf{Date:}\quad 28/04/2026 & \\
\end{tabular}
\end{center}
\hrule\vspace{1.5em}

\subsection*{Overview}
This lecture has three conceptual parts.  In the first part we pick up directly
from the covariant electromagnetic action of Lecture~06, choose coordinate time
as the worldline parameter, and extract the explicit three-dimensional Lagrangian.
We identify the \textbf{scalar potential} $\phi$ and the \textbf{vector potential}
$\Avec$ as the temporal and spatial parts of the covariant four-potential $A_\mu$,
and carefully distinguish the \textbf{canonical momentum} from the \textbf{mechanical
momentum}.

In the second part we apply the Euler--Lagrange equations to that Lagrangian.
The key technical step is evaluating the total time derivative of the potential
along the particle's worldline via the chain rule; a careful bookkeeping of the
Levi-Civita identity then causes two gradient terms to cancel and the familiar
\textbf{Lorentz force law} drops out.

In the third part we identify a profound symmetry of the entire formulation:
\textbf{gauge invariance}.  Shifting $A_\mu$ by the four-gradient of any smooth
scalar function $f$ changes the action by a pure boundary term, leaving the
equations of motion and the physical fields $\bE$ and $\bB$ unchanged.  We
close with a manifestly covariant re-derivation of the equations of motion that
uses the \textbf{proper time} $\tau$ as the parameter; the calculation keeps
Lorentz invariance explicit at every step and delivers the electromagnetic
\textbf{field tensor} $F^{\mu\nu}$.

% ==============================================================================
\subsection{From the Covariant Action to the Explicit Lagrangian}\label{sec:lec07:lagrangian}
% ==============================================================================

\subsubsection*{Physical Motivation}

In Lecture~06 we derived the total action for a relativistic charged particle
in an external electromagnetic field:
\begin{equation}\label{eq:lec07:total-action}
    S = -mc\int ds - \frac{e}{c}\int A_\mu\,dx^\mu .
\end{equation}
This expression is manifestly Lorentz invariant, but it is written without
committing to any particular parameter along the worldline.  To derive explicit
equations of motion we must \emph{choose a parameter}.  The natural choice is
coordinate time $t$: this breaks 4-covariance at the level of the Lagrangian
but is perfectly consistent and reproduces Newton's equation in the
non-relativistic limit.  The translation into $L(t)$ is a straightforward
but instructive exercise in decomposing four-vectors.

\subsubsection*{Decomposing the Four-Potential}

The electromagnetic four-potential with a raised (contravariant) index is
defined as
\begin{equation}\label{eq:lec07:fourpot-upper}
    A^\mu = (\phi,\,\Avec),
\end{equation}
where $\phi$ is the \textbf{scalar potential} and $\Avec$ the
\textbf{three-vector potential}.  Lowering with the Minkowski metric
$\eta_{\mu\nu} = \mathrm{diag}(+1,-1,-1,-1)$:
\begin{equation}\label{eq:lec07:fourpot-lower}
    A_\mu = \eta_{\mu\nu}A^\nu
    = (\phi,\,-\Avec).
\end{equation}
The time component is preserved ($A_0 = A^0 = \phi$), while the spatial
components pick up a minus sign: $A_i = -A^i$.  This sign is a direct
consequence of the $(+,-,-,-)$ signature; it appears in every contraction
involving spatial indices.

\subsubsection*{Expanding the Interaction Term}

With $x^\mu(t) = (ct, \bvec{x}(t))$ as the parameterisation, the coordinate
differentials are $dx^\mu = (c\,dt,\,d\bvec{x})$.  Contracting with $A_\mu$:
\begin{align}
    A_\mu\,dx^\mu
    &= A_0\,dx^0 + A_i\,dx^i
     = \phi\cdot c\,dt + (-A^i)\,dx^i
     = \phi\,c\,dt - \Avec\cdot d\bvec{x} \notag \\
    &= \left(c\phi - \Avec\cdot\bv\right)dt ,
    \label{eq:lec07:Amu-dxmu}
\end{align}
where $\bv = d\bvec{x}/dt$ is the three-velocity.  The interaction term in
the action becomes:
\begin{equation}\label{eq:lec07:Sint}
    -\frac{e}{c}\int A_\mu\,dx^\mu
    = -\frac{e}{c}\int\!\left(c\phi - \Avec\cdot\bv\right)dt
    = \int\!\left(-e\phi + \frac{e}{c}\,\Avec\cdot\bv\right)dt .
\end{equation}
Adding the free-particle Lagrangian derived in Lecture~06
($L_{\text{free}} = -mc^2\!\sqrt{1-v^2/c^2}$), the complete Lagrangian is:
\begin{equation}\label{eq:lec07:lagrangian}
    \boxed{L = -mc^2\!\sqrt{1-\frac{v^2}{c^2}} - e\phi + \frac{e}{c}\,\Avec\cdot\bv .}
\end{equation}

\textit{Physical significance:} The three terms encode, in order: the
relativistic inertia (whose non-relativistic expansion yields $\tfrac{1}{2}mv^2$
plus the irrelevant constant $-mc^2$); the electrostatic potential energy $e\phi$;
and a \emph{velocity-dependent} magnetic coupling $\tfrac{e}{c}\Avec\cdot\bv$.
The magnetic coupling is the origin of the $\bvec{v}\times\bB$ term in the
Lorentz force: it is the price paid for the field carrying four-momentum.

% ==============================================================================
\subsection{Equations of Motion: Deriving the Lorentz Force}\label{sec:lec07:EOM}
% ==============================================================================

\subsubsection*{Physical Motivation}

Because $\phi$ and $\Avec$ depend on position, the action is not invariant
under spatial translations, and the canonical momentum is not conserved.  The
task is to apply the Euler--Lagrange equations
\begin{equation}\label{eq:lec07:EL}
    \frac{d}{dt}\!\left(\frac{\partial L}{\partial v^i}\right) - \frac{\partial L}{\partial x^i} = 0
\end{equation}
and work out both sides explicitly.  The key subtlety is the total time
derivative of the potential along the moving worldline, which requires the
chain rule.  A vector identity (proved via Levi-Civita symbols) then causes
two gradient terms to cancel, leaving a clean result.

\subsubsection*{The Canonical Momentum}

Differentiating $L$ with respect to $v^i$:
\begin{equation}\label{eq:lec07:canonical-momentum}
    p_i^{\text{can}} \;\equiv\; \frac{\partial L}{\partial v^i}
    = \underbrace{\frac{mv^i}{\sqrt{1-v^2/c^2}}}_{\displaystyle p_i^{\text{mech}}}
      + \frac{e}{c}\,A^i .
\end{equation}
Two contributions appear.  The \textbf{mechanical momentum}
$p_i^{\text{mech}} = \gamma mv^i$ is the same quantity derived from the free
Lagrangian in Lecture~06.  The field-dependent correction $\tfrac{e}{c}A^i$
is new — it arises from the velocity-dependent coupling and makes the canonical
momentum \textbf{gauge dependent} (since $A^i$ changes under a gauge
transformation, as we will see).

\begin{remark}
In quantum mechanics (and quantum field theory), canonical commutation
relations are imposed between position $x^i$ and the \emph{canonical} momentum
$p_j^{\text{can}}$.  The minimal-coupling substitution
$\bvec{p}^{\text{can}}\to -i\hbar\nabla$ therefore replaces the
\emph{mechanical} momentum by $-i\hbar\nabla - \tfrac{e}{c}\Avec$.  This is
the origin of the covariant derivative in quantum electrodynamics and is
a non-trivial physical consequence of the distinction between the two momenta.
\end{remark}

\subsubsection*{The Generalised Force: A Vector Identity via Levi-Civita Symbols}

The right-hand side of the Euler--Lagrange equation requires $\partial L/\partial x^i$
with $\bv$ held fixed:
\begin{equation}\label{eq:lec07:dLdx-raw}
    \frac{\partial L}{\partial x^i}
    = -e\,\partial_i\phi + \frac{e}{c}\,\frac{\partial(\Avec\cdot\bv)}{\partial x^i}
    = -e\,\partial_i\phi + \frac{e}{c}\,v^j\,\partial_i A^j .
\end{equation}
(Since $\bv$ does not depend on $\bvec{x}$, only the spatial components of
$\Avec$ are differentiated.)

We now rewrite $v^j\,\partial_i A^j$ using the Levi-Civita symbol.  Recall
the identity
\begin{equation}\label{eq:lec07:levi-civita-id}
    \epsilon_{ijk}\,\epsilon_{klm} = \delta_{il}\,\delta_{jm} - \delta_{im}\,\delta_{jl} .
\end{equation}
The $i$-th component of $\bvec{v}\times(\nabla\times\Avec)$ is
\begin{align}
    \left[\bvec{v}\times(\nabla\times\Avec)\right]^i
    &= \epsilon^{ijk}\,v^j\,\left(\nabla\times\Avec\right)^k
     = \epsilon^{ijk}\,v^j\,\epsilon_{klm}\,\partial^l A^m \notag \\
    &= \left(\delta^i_l\,\delta^j_m - \delta^i_m\,\delta^j_l\right)v^j\,\partial^l A^m
     = v^j\,\partial^i A^j - v^j\,\partial^j A^i .
    \label{eq:lec07:cross-expand}
\end{align}
With $\bv$ treated as spatially constant:
$v^j\,\partial^i A^j = \partial^i(v^j A^j) = \partial_i(\bv\cdot\Avec)$ (using
$\partial^i = -\partial_i$ for spatial indices with our metric, but $\partial_i(\bv\cdot\Avec)$
suffices here since $\bv$ is constant).  Therefore:
\begin{equation}\label{eq:lec07:cross-identity}
    v^j\,\partial_i A^j
    = \left[\bvec{v}\times(\nabla\times\Avec)\right]^i + v^j\,\partial_j A^i
    = \left[\bvec{v}\times\bB\right]^i + (\bv\cdot\nabla)A^i ,
\end{equation}
where $\bB \equiv \nabla\times\Avec$.  Substituting back:
\begin{equation}\label{eq:lec07:dLdx-final}
    \frac{\partial L}{\partial x^i}
    = -e\,\partial_i\phi + \frac{e}{c}\,(\bv\cdot\nabla)A^i + \frac{e}{c}\,[\bvec{v}\times\bB]^i .
\end{equation}

\subsubsection*{Total Time Derivative Along the Worldline}

The left-hand side of the Euler--Lagrange equation involves
$\frac{d}{dt}(\gamma mv^i + \tfrac{e}{c}A^i)$.  The potential
$\bvec{A} = \Avec(\bvec{x}(t),t)$ depends on both position and time; as the
particle moves, both change.  Applying the chain rule:
\begin{equation}\label{eq:lec07:chain-rule}
    \frac{dA^i}{dt}
    = \underbrace{\frac{\partial A^i}{\partial t}}_{\text{field changes at fixed point}}
      +\; \underbrace{v^j\,\frac{\partial A^i}{\partial x^j}}_{\text{particle moves to new point}}
    = \frac{\partial A^i}{\partial t} + (\bv\cdot\nabla)A^i .
\end{equation}
The first term is the explicit time dependence of the field; the second arises
because the particle drifts to a location where the field has a different value.

\subsubsection*{Assembling the Euler--Lagrange Equation}

Inserting~\eqref{eq:lec07:chain-rule} and~\eqref{eq:lec07:dLdx-final} into the
Euler--Lagrange equation:
\begin{equation}\label{eq:lec07:EL-assembled}
    \frac{d(\gamma mv^i)}{dt}
    + \frac{e}{c}\!\left(\frac{\partial A^i}{\partial t}
      + \cancel{(\bv\cdot\nabla)A^i}\right)
    = -e\,\partial_i\phi + \frac{e}{c}\,\cancel{(\bv\cdot\nabla)A^i}
      + \frac{e}{c}\,[\bvec{v}\times\bB]^i .
\end{equation}
The two $(\bv\cdot\nabla)A^i$ terms cancel identically.  Rearranging:
\begin{equation}\label{eq:lec07:EL-simplified}
    \frac{d(\gamma mv^i)}{dt}
    = -e\,\partial_i\phi - \frac{e}{c}\,\frac{\partial A^i}{\partial t}
      + \frac{e}{c}\,[\bvec{v}\times\bB]^i .
\end{equation}

\subsubsection*{Definitions of the Physical Fields and the Lorentz Force}

Recognising the two combinations that appear, we define:
\begin{align}
    \bE &\;\equiv\; -\nabla\phi - \frac{1}{c}\frac{\partial\Avec}{\partial t} ,
    \label{eq:lec07:E-def}\\
    \bB &\;\equiv\; \nabla\times\Avec .
    \label{eq:lec07:B-def}
\end{align}
Substituting into~\eqref{eq:lec07:EL-simplified}:
\begin{equation}\label{eq:lec07:lorentz-force}
    \boxed{\frac{d\bvec{p}^{\,\text{mech}}}{dt}
    = e\!\left(\bE + \frac{\bv}{c}\times\bB\right).}
\end{equation}
This is the \textbf{Lorentz force law}.

\textit{Physical significance:}  We have derived—purely from symmetry (Lorentz
invariance) and the single coupling $-\tfrac{e}{c}\int A_\mu\,dx^\mu$—the
complete electromagnetic force on a charged relativistic particle.  The form was
not assumed; it followed from the Euler--Lagrange equations applied to the
action.  The cancellation of the two $(\bv\cdot\nabla)A^i$ terms is not
accidental: it is mandated by gauge invariance, as will become clear in
Section~\ref{sec:lec07:gauge}.

\subsubsection*{Why the Equation of Motion Is Lorentz Invariant}

The equation~\eqref{eq:lec07:lorentz-force} looks correct but its Lorentz
invariance is not manifest: $\bvec{v}$ is not a four-vector, $\bE$ and $\bB$
transform into each other under boosts, and it is laborious to verify explicitly
that both sides transform the same way.

The key insight is: because we \emph{derived} the equation from the
Lorentz-invariant action~\eqref{eq:lec07:total-action}, it \emph{must} be
Lorentz invariant.  Any Lorentz transformation merely relabels which inertial
frame we call ``home''; since the action is the same in all frames, so are its
Euler--Lagrange equations, up to a change of variables.  This is a powerful
application of the symmetry-first approach: we enforce invariance at the level
of the action (one scalar object) and the equations of motion inherit it
automatically, even if their form looks frame-dependent.

% ==============================================================================
\subsection{Gauge Invariance}\label{sec:lec07:gauge}
% ==============================================================================

\subsubsection*{Physical Motivation}

The Lagrangian~\eqref{eq:lec07:lagrangian} is written in terms of $\phi$ and
$\Avec$ rather than the directly measurable fields $\bE$ and $\bB$.  This
raises a natural question: how much freedom do we have in choosing $A_\mu$
while describing the \emph{same} physical situation?  Answering this question
reveals a deep structure: the electromagnetic field is not uniquely described
by a single $A_\mu$, but by an entire \emph{equivalence class} of potentials.
The freedom to choose among them is called \textbf{gauge freedom}, and it is
one of the defining features of every fundamental force in nature.

\subsubsection*{The Gauge Transformation}

A \textbf{gauge transformation} replaces the four-potential by
\begin{equation}\label{eq:lec07:gauge-transform}
    \boxed{A_\mu \;\longrightarrow\; A_\mu' = A_\mu + \partial_\mu f ,}
\end{equation}
where $f = f(x^\alpha)$ is any smooth scalar function of spacetime.  In
components, using $\partial_\mu f = (\tfrac{1}{c}\partial_t f,\,\nabla f)$:
\begin{align}
    \phi  &\;\longrightarrow\; \phi + \frac{1}{c}\frac{\partial f}{\partial t} ,
    \label{eq:lec07:gauge-phi}\\
    \Avec &\;\longrightarrow\; \Avec + \nabla f .
    \label{eq:lec07:gauge-A}
\end{align}
The function $f$ is completely arbitrary: there are infinitely many gauge-related
descriptions of any given electromagnetic field.

\subsubsection*{The Action Changes by a Boundary Term Only}

Under~\eqref{eq:lec07:gauge-transform}, the action changes by
\begin{align}
    \delta S
    &= -\frac{e}{c}\int(\partial_\mu f)\,dx^\mu
     = -\frac{e}{c}\int \frac{\partial f}{\partial x^\mu}\,dx^\mu
     = -\frac{e}{c}\int df
     = -\frac{e}{c}\bigl[f(x_{\text{final}}) - f(x_{\text{initial}})\bigr].
    \label{eq:lec07:delta-S}
\end{align}
The integral of an exact differential $df$ telescopes to the endpoints.  Since
the variational principle holds the endpoints fixed ($\delta x^\mu = 0$ there),
$\delta S$ is a constant that depends only on the fixed boundary values of $f$
and is therefore \textbf{irrelevant for the equations of motion}.

\textit{Physical significance:}  The gauge freedom means that $A_\mu$ is not a
physical observable.  Any two potentials related by
$A_\mu' = A_\mu + \partial_\mu f$ represent the \emph{same} physical state.
The actual physical content is encoded in the equivalence class of all such
potentials.  Singling out one representative — a \textbf{gauge choice} — is
analogous to choosing a coordinate system: it simplifies calculations but
carries no physical meaning.  Different gauges that are useful in different
contexts include the \emph{Lorenz gauge} ($\partial^\mu A_\mu = 0$, useful for
manifestly covariant calculations) and the \emph{Coulomb gauge} ($\nabla\cdot\Avec = 0$,
useful for non-relativistic problems).

\subsubsection*{Uniqueness of the Action up to Boundary Terms}

This observation also answers a deeper question: \emph{how unique is the action?}
Two actions that yield identical equations of motion must have identical
functional derivatives (``functional gradients'').  By the same argument that
shows two ordinary functions with equal gradients differ by at most a constant,
two actions with identical Euler--Lagrange equations differ by at most a
functional of the boundary values — i.e.\ a total time derivative.  Adding a
total derivative to $L$ (equivalently, a boundary term to $S$) never changes
the equations of motion.  The gauge transformation is precisely an instance of
this: it adds $-\tfrac{e}{c}df$ to the Lagrangian, which is a total derivative.

\subsubsection*{The Physical Fields Are Gauge Invariant}

We verify directly that $\bE$ and $\bB$ are unchanged under~\eqref{eq:lec07:gauge-phi}--\eqref{eq:lec07:gauge-A}:
\begin{align}
    \bE' &= -\nabla\!\left(\phi + \frac{\partial_t f}{c}\right)
             - \frac{1}{c}\frac{\partial}{\partial t}\!\left(\Avec + \nabla f\right)
           = \bE
             - \frac{1}{c}\nabla\partial_t f
             + \frac{1}{c}\partial_t\nabla f = \bE ,
    \label{eq:lec07:E-invariant}\\
    \bB' &= \nabla\times(\Avec + \nabla f)
           = \bB + \underbrace{\nabla\times\nabla f}_{= 0} = \bB .
    \label{eq:lec07:B-invariant}
\end{align}
In~\eqref{eq:lec07:E-invariant} the two extra terms cancel because partial
derivatives commute; in~\eqref{eq:lec07:B-invariant} the curl of a gradient is
identically zero.

Both physical fields are gauge invariant: $\bE' = \bE$ and $\bB' = \bB$.  This
is consistent with the Lorentz force law~\eqref{eq:lec07:lorentz-force}, which
is written entirely in terms of $\bE$ and $\bB$, not $\phi$ or $\Avec$.

% ==============================================================================
\subsection{Manifestly Covariant Derivation of the Equations of Motion}\label{sec:lec07:covariant}
% ==============================================================================

\subsubsection*{Physical Motivation}

In Section~\ref{sec:lec07:EOM} we derived the Lorentz force using coordinate
time $t$, which breaks the democratic treatment of space and time.  The result
is correct, but its Lorentz invariance is hidden.  There is a conceptually
cleaner approach: work with the \textbf{proper time} $\tau$ as the parameter,
never commit to a specific frame, and vary the covariant
action~\eqref{eq:lec07:total-action} directly.  The payoff is a single,
manifestly covariant equation of motion in which the Lorentz symmetry is
self-evident, and it naturally introduces the \textbf{electromagnetic field
tensor} $F^{\mu\nu}$ — a single, gauge-invariant, Lorentz-covariant object that
encodes both $\bE$ and $\bB$.

\subsubsection*{Setup: Proper Time as Parameter}

Since $ds = c\,d\tau$, the action becomes
\begin{equation}\label{eq:lec07:action-proper}
    S = -mc^2\int d\tau - \frac{e}{c}\int A_\mu\,u^\mu\,d\tau ,
\end{equation}
where the \textbf{four-velocity} $u^\mu = dx^\mu/d\tau = \gamma(c,\bv)$
satisfies $u^\mu u_\mu = c^2$.

\subsubsection*{Variation of the Worldline}

We shift the worldline: $x^\mu(\tau) \to x^\mu(\tau) + \delta x^\mu(\tau)$,
with $\delta x^\mu = 0$ at both endpoints.  The resulting changes are:
\begin{align}
    \delta u^\mu &= \frac{d(\delta x^\mu)}{d\tau} ,
    \label{eq:lec07:delta-u}\\
    \delta A_\mu &= (\partial_\nu A_\mu)\,\delta x^\nu
    \quad\text{(Taylor expansion to first order).}
    \label{eq:lec07:delta-A}
\end{align}

\subsubsection*{First Variation of the Free-Particle Term}

From $c^2 d\tau^2 = \eta_{\mu\nu}dx^\mu dx^\nu$ one finds
$c^2\,\delta(d\tau) = u_\nu\,d(\delta x^\nu)$, so:
\begin{equation}\label{eq:lec07:delta-ds}
    -mc^2\,\delta\!\int d\tau
    = -m\!\int u_\nu\,\frac{d(\delta x^\nu)}{d\tau}\,d\tau
    \overset{\text{IBP}}{=}
    +m\!\int \frac{du_\nu}{d\tau}\,\delta x^\nu\,d\tau ,
\end{equation}
where integration by parts was used and the boundary terms vanish because
$\delta x^\nu = 0$ at the endpoints.

\subsubsection*{First Variation of the Interaction Term}

\begin{align}
    -\frac{e}{c}\,\delta\!\int A_\mu\,u^\mu\,d\tau
    &= -\frac{e}{c}\!\int\!
       \Bigl[(\partial_\nu A_\mu)\,\delta x^\nu\,u^\mu
       + A_\mu\,\frac{d(\delta x^\mu)}{d\tau}\Bigr]\,d\tau .
    \label{eq:lec07:delta-interaction}
\end{align}
Integrating the second piece by parts (boundary terms vanish):
\begin{equation}\label{eq:lec07:IBP-interaction}
    -\frac{e}{c}\!\int A_\mu\,\frac{d(\delta x^\mu)}{d\tau}\,d\tau
    = +\frac{e}{c}\!\int \frac{dA_\mu}{d\tau}\,\delta x^\mu\,d\tau
    = +\frac{e}{c}\!\int (\partial_\nu A_\mu)\,u^\nu\,\delta x^\mu\,d\tau .
\end{equation}
Combining~\eqref{eq:lec07:delta-interaction} and~\eqref{eq:lec07:IBP-interaction},
and renaming the dummy index in the first term ($\mu \leftrightarrow \nu$) so
that $\delta x^\mu$ factors out:
\begin{equation}\label{eq:lec07:interaction-combined}
    -\frac{e}{c}\,\delta\!\int A_\mu\,u^\mu\,d\tau
    = \frac{e}{c}\!\int
      \bigl[(\partial_\mu A_\nu - \partial_\nu A_\mu)\,u^\nu\bigr]\,\delta x^\mu\,d\tau .
\end{equation}

\subsubsection*{The Field Tensor and the Covariant Equation of Motion}

Adding~\eqref{eq:lec07:delta-ds} and~\eqref{eq:lec07:interaction-combined} and
setting $\delta S = 0$ for arbitrary $\delta x^\mu$:
\begin{equation}\label{eq:lec07:variation-zero}
    \int\!\left[
        m\,\frac{du_\mu}{d\tau}
        - \frac{e}{c}\,(\partial_\mu A_\nu - \partial_\nu A_\mu)\,u^\nu
    \right]\delta x^\mu\,d\tau = 0 .
\end{equation}
Since $\delta x^\mu$ is an arbitrary function, the integrand must vanish at
every point.  We define the \textbf{electromagnetic field tensor}:
\begin{equation}\label{eq:lec07:Fmunu-def}
    \boxed{F_{\mu\nu} \;\equiv\; \partial_\mu A_\nu - \partial_\nu A_\mu ,}
\end{equation}
an antisymmetric rank-2 covariant tensor.  Raising both indices (or,
equivalently, using $F^{\mu\nu} = \partial^\mu A^\nu - \partial^\nu A^\mu$),
the equation of motion takes the manifestly covariant form:
\begin{equation}\label{eq:lec07:covariant-EOM}
    \boxed{m\,\frac{du^\mu}{d\tau}
    = \frac{e}{c}\,F^{\mu}_{\ \nu}\,u^\nu .}
\end{equation}

\textit{Physical significance:}  This is the \textbf{relativistic Lorentz force
law} in fully covariant form.  It is a single 4-component equation: the spatial
components reproduce~\eqref{eq:lec07:lorentz-force}, while the time component
gives the power equation $dE/dt = e\bE\cdot\bv$ (the rate at which the field does
work on the particle).  Lorentz invariance is manifest: $F^{\mu}_{\ \nu}$ is a
tensor and $u^\nu$ is a four-vector, so both sides transform identically under
any Lorentz transformation.  There is no need to verify invariance term by term.

\subsubsection*{Structure of the Field Tensor: Encoding $\bE$ and $\bB$}

The tensor $F^{\mu\nu} = -F^{\nu\mu}$ has $\binom{4}{2} = 6$ independent
components.  Comparing the definition~\eqref{eq:lec07:Fmunu-def} with the
definitions~\eqref{eq:lec07:E-def}--\eqref{eq:lec07:B-def}:
\begin{itemize}
    \item \textbf{Time-space components:} $F^{0i}$ encodes the electric field.
    \item \textbf{Space-space components:} $F^{ij}$ encodes the magnetic field
      via $F^{ij} = -\epsilon^{ijk}B_k$.
\end{itemize}
All six physical field components ($\bE$ and $\bB$) are packaged into a single
Lorentz tensor.  Under a Lorentz boost, $F^{\mu\nu}$ transforms as a tensor
and mixes the electric and magnetic components — which is precisely the
well-known fact that a purely electric field in one frame contains a magnetic
component in a boosted frame.

\subsubsection*{Gauge Invariance of the Field Tensor}

Under $A_\mu \to A_\mu + \partial_\mu f$:
\begin{equation}\label{eq:lec07:F-gauge-inv}
    F_{\mu\nu} \;\to\; F_{\mu\nu}
    + \underbrace{(\partial_\mu\partial_\nu f - \partial_\nu\partial_\mu f)}_{= 0}
    = F_{\mu\nu} .
\end{equation}
Partial derivatives commute, so $F_{\mu\nu}$ is \textbf{identically gauge
invariant}.  This is the deeper reason why $\bE$ and $\bB$ — which are the
components of $F^{\mu\nu}$ — are gauge invariant, as we verified directly in
Section~\ref{sec:lec07:gauge}.

% ==============================================================================
\subsection*{Summary and Looking Ahead}
% ==============================================================================

\begin{enumerate}
    \item \textbf{Explicit Lagrangian:}
      $L = -mc^2\sqrt{1-v^2/c^2} - e\phi + \tfrac{e}{c}\Avec\cdot\bv$.
      The three terms are inertia, electrostatic coupling, and magnetic coupling.

    \item \textbf{Canonical vs.\ mechanical momentum:}
      $\bvec{p}^{\text{can}} = \gamma m\bv + \tfrac{e}{c}\Avec$.
      The canonical momentum enters Poisson brackets and canonical quantisation;
      the mechanical momentum appears in the force law.

    \item \textbf{Lorentz force derived from symmetry:}
      The Euler--Lagrange equations, combined with the chain rule for
      $dA^i/dt$ along the worldline and a Levi-Civita identity, yield
      $d\bvec{p}^{\text{mech}}/dt = e(\bE + \bvec{v}/c\times\bB)$.
      The two $(\bv\cdot\nabla)A^i$ terms cancel.

    \item \textbf{Gauge invariance:} $A_\mu \to A_\mu + \partial_\mu f$ changes
      the action by a boundary term only.  Physical states correspond to
      equivalence classes of potentials.  The fields $\bE$ and $\bB$ are gauge
      invariant.  Different gauge choices simplify different calculations.

    \item \textbf{Uniqueness of the action:} Two Lagrangians with identical
      equations of motion differ by at most a total time derivative.  The gauge
      transformation generates precisely such a difference.

    \item \textbf{Covariant equation of motion:}
      $m\,du^\mu/d\tau = \tfrac{e}{c}F^{\mu}_{\ \nu}u^\nu$,
      derived by varying the covariant action with proper time as parameter
      and integrating by parts.

    \item \textbf{The field tensor:} $F_{\mu\nu} = \partial_\mu A_\nu - \partial_\nu A_\mu$
      is antisymmetric, gauge invariant, and packs all six electromagnetic
      field components ($\bE$ and $\bB$) into a single Lorentz tensor.

\end{enumerate}

In Lecture~08 we will give the field tensor $F^{\mu\nu}$ a life of its own:
we will write the \textbf{action for the electromagnetic field itself} (not
just for a particle in the field) and show that its Euler--Lagrange equations
reproduce \textbf{Maxwell's equations} in covariant form.  We will also meet
the \textbf{energy-momentum tensor} $\emtensor$ of the electromagnetic field
and connect it to conservation laws via Noether's theorem.
